\documentclass{article}
\usepackage{graphicx} % Required for inserting images
\usepackage{gensymb}
\usepackage[round, authoryear]{natbib} 
\bibliographystyle{plainnat} 


\title{DS9600 Final Project \\ \textbf {Beyond the Black Box: A Multimodal XAI Framework for Lumbar Spinal Stenosis Grading via Integrated MRI Analysis and LLM-Driven Clinical Reasoning}}
\author{Ahmed Mohamed, (Helen) Yixiu He, Yan Yi Li, Tim Nguyen}
\date{March 2026}

\begin{document}
\maketitle


\tableofcontents

\clearpage

\section{Problem Statement}
\label{problem_statement}
In the field of medical imaging, medical practitioners deal with a rising diagnostic workload.  To mitigate this issue, many choose to extend work hours and increase productivity expectations \citep{Jing2025}. However, this promotes job dissatisfaction and burnout. In private practice, the burnout rate is 46\% for radiologists. This result strongly correlates with practitioners who took evenings, overnight and weekends calls \citep{Parikh2023}. Not just in private sector, but burn out also happens in acadmedia, where burnout overall prevalence rate is 37.4\% \citep{Higgins2021}. Even though practices provide solution to mitigate this issue such as flexible scheduling, part-time and remote work options, teleradiology services \citep{Rawson2024}, these only address the surface problems instead of targeting the root cause, which is the shortage of radiologists  \citep{Jing2025}. With the increasing number of patients but shortage of professionals to diagnose the images, it causes huge delay and decrease in quality in health care providing. Due to an increasing of already heavy work flow, combining with ineffective operation and shortage of radiologists, this leads to excessive workload, with visual and mental fatigue. These can lead to imaging delays, causing delays of diagnosis and treatment \citep{Cournane2016}. In 2025, a survey reports that 11.4\% of patient had to stop working while waiting for a result of a scan. The wait time can vary between one up to three months \citep{nanos}. 
For that issue, we are interested in building a pipeline to address Lumbar Spinal Stenosis (LSS), but the general idea is applicable for any medical imaging. In order to address this problem, researchers developed a pipeline to utilize the power of computer vision to aid medical practitioners in diagnosing imaging. In particular, there exists a pipeline to use U-Net architecture to perform ROI segmentation \citep{Silveira2025}, ConvNet in image classification and Grad-CAM to generate feature importance heatmap \citep{Selvaraju_2019}. However, in a clinical survey, 56\% radiologists mentioned a lack of transparency since this pipeline utilizes black box models \citep{aravazhi2024}, there is lack of motivation behind predictions. Furthermore, they state that they prefer textual explanation \citep{gale2018producingradiologistqualityreportsinterpretable, kayser2024fooloncecontrastingtextual} , which our model comes to play. Here, we propose the use of multimodal large language model to attempt to explain the motivation behind prediction using ROI and heatmap. 

\section{Rationale and Background}
\subsection{Clinical Burden and Diagnostic Variability}
Lumbar Spinal Stenosis (LSS) is a leading cause of neurogenic claudication and a primary indication for spinal surgery in adults over 60, affecting an estimated 11\% of the general population \citep{Munakomi2023}. Magnetic Resonance Imaging (MRI) is the gold standard for diagnosis; however, the manual assessment of multiple spinal levels (L1–S1) is labor-intensive and suffers from significant inter-observer variability. Studies indicate that even expert radiologists show only “moderate” agreement (Cohen’s kappa 0.40–0.57) when grading stenosis severity \citep{AlTameemi2017, Holc2023}. This diagnostic inconsistency can lead to misaligned treatment plans, varying from unnecessary surgeries to delayed interventions. Due to these two reasons, developing a tool to support diagnosis process and improve accuracy is essential for this condition. 

\subsection{The Rise and Limitations of Deep Learning in Biology}
Recent advancements in Deep Learning (DL) have demonstrated that Convolutional Neural Networks (CNNs) can achieve expert-level accuracy in LSS classification. Models utilizing architectures like U-Net for segmentation with CNN-based classifiers have reached accuracy exceeding 0.94 \citep{Liu2025}. Despite this high performance, these models primarily function as “black boxes.” In a clinical survey by Bioinformation, 56\% of radiologists cited a lack of transparency as a major barrier to AI adoption \citep{aravazhi2024}. While visual tools like Grad-CAM (Gradient-weighted Class Activation Mapping) highlight “hotspots” of model attention, they are often criticized in medical literature for being “noisy” and anatomically imprecise, failing to distinguish between pathologically relevant areas and background noise \citep{Houssein2025, Muhammad2024}.

\subsection{The Unfeasiability of Interpretable AI}
One option for this problem is adopting interpretable models instead. In high-stake industries such as medicine and law, the use of interpretable models such as linear regression, logistic regression and decision list is favoured due to their transparency, as opposed to black box models where the content is opaque even to its developers \citep{wang2019hybridpredictivemodelinterpretable}. Many literatures encourage the pursuit for simple models allowing the understanding of motivation behind decisions, giving the stakeholders the necessary reasonings for decision making \citep{rudin2019stopexplainingblackbox}. However, also according to \cite{rudin2019stopexplainingblackbox}, it is particularly challenging, if not impossible, to utilize these models for complex tasks, in which one of them is computer vision and imaging. In this pipeline, since we exploit the power of computer vision techniques as decision making assistants, it is virtually not possible to implement simple models that can understand images in an effective manner. That leaves us one other possible solution for this conundrum.

\subsection{Explainable AI \& its “Semantic Gap”}

In this proposal, we suggest the use of explainable AI (XAI) to mitigate the opaqueness of black box models. In the current pipeline, a heatmap generated by using Grad-CAM to highlight the important area on the imaging that influences the prediction. However, there is a disconnection between pixel-level importance and clinical reasoning. A clinician does not make a diagnosis based on a heatmap; they make it based on anatomical observations such as “thecal sac compression” or “posterior disc bulging.” Research in human-computer interaction (HCI) suggests that clinicians strongly prefer textual explanations that mirror the structure of a radiology report over purely visual saliency maps (Gale et al., 2018; Kayser et al., 2024). Furthermore, visual-only explanations can lead to automation bias, where clinicians may trust a correct prediction for the wrong underlying reason (e.g., the model focused on a surgical clip rather than the spinal canal) \citep{Houssein2025}. This problem also aligns with a study by \cite{panboonyuen2026seeingisntbelievinganalysis}, where they found heatmaps generated by Grad-CAM fail to represent true motivation behind a decision. Therefore, to make heatmap more accessible and helpful to medical practitioners, textual explanation is needed, hence the adoption of Large Language Models (LLMs). 

\subsection{Identifying the Gap: The Need for Vision-Language Fusion}
Large Language Models (LLMs) are models that are trained on a large dataset of text. They are able to provide concise and easy to comprehend textual explanations. In this proposal, we are suggesting the use of Multimodal Large Language Models (MLLMs). MLLMs are trained to connect different modalities such as visual, audio and text together to provide a comprehensive presentation that connect the three sectors. Here, we are proposing the use of MR imaging, MRI segmented ROI, heatmap generated by Grad-CAM and prediction by the pipeline as inputs for the MLLMs to generate textual explanation.  
However, while LLMs have shown remarkable ability in synthesizing medical information, they lack “grounding” in raw imaging data when used in isolation, leading to “hallucinations” \citep{pal2024geminigoesmedschool, Omar2025, Yang2025} . The identified gap is the lack of a pipeline that translates quantitative, segmented anatomical features and visual attention patterns into structured, clinical-grade natural language. 
By anchoring an LLM’s reasoning in morphometric data (e.g., exact canal area measurements) and spatial overlap metrics (how much the AI’s attention overlaps with the thecal sac), our proposal creates a novel, interdisciplinary “translation layer.” This addresses the technical challenge of XAI while fulfilling the clinical requirement for evidence-based diagnostic support, ultimately moving AI from a “scoring tool” to a “transparent diagnostic partner.”

\section{Proposed ML/AI System or Solution}
This study aims to develop an interpretable deep learning framework for automated severity assessment of Lumbar spinal stenosis (LSS) using axial T2-weighted Magnetic Resonance Imaging (MRI), accompanied by a structured text-based clinical explanation of the model’s decision. The proposed framework consists of two sequential components: anatomical structure segmentation with quantitative morphometric analysis, followed by stenosis severity classification with integrated explainability.

\subsection{Datasets}
For anatomical segmentation, we will use a publicly available lumbar spine MRI dataset from Mendeley Data (Dataset ID: k57fr854j2), derived from an anonymized clinical cohort of patients with lower back pain \citep{lumbardataset}. The dataset includes T2-weighted lumbar MRI scans and a companion annotated subset (Dataset ID: k57fr854j2.1) with expert-generated segmentation masks for key anatomical structures (e.g., intervertebral discs and spinal canal) \citep{labelImg}, which will serve as the reference standard for model training and evaluation. A curated subset (around 500 mid-sagittal slices) will be used. Also, the dataset will be split at the patient level into training/validation (80\%) and testing (20\%) sets to prevent data leakage. For stenosis severity classification, we will utilize the RSNA Lumbar Spine Degenerative Classification Dataset (n = 2,697), which provides vertebral-level stenosis grades (normal, mild, moderate, severe) \citep{Richards2026}. The segmentation model trained on the Mendeley dataset will be applied to the RSNA dataset to extract consistent anatomical regions of interest (ROIs) and quantitative features across all cases, enabling standardized input for downstream classification.

\paragraph{Phase 1: Quantitative ML Prediction Framework}

\clearpage

\bibliography{references.bib}

\end{document}
